\documentclass[12pt]{article}
\usepackage[utf8]{inputenc}
\usepackage[T1]{fontenc}
\usepackage[margin=1in]{geometry}
\usepackage{setspace}
\usepackage{xcolor}
\usepackage[colorlinks=true, urlcolor=blue!60!black, linkcolor=blue!60!black]{hyperref}
\setlength{\parindent}{2em}
\setlength{\parskip}{1.25em}
\renewcommand{\baselinestretch}{1.0}
\begin{document}
\begin{center}
{\LARGE \textit{Moore v. Harper}, 600 U.S. 1 (2023)\par}
\vspace{0.8em}
{\large Case Brief\par}
\end{center}
\vspace{1em}

\section{Memory Jogger}
The Elections Clause does not insulate state legislatures from ordinary state-court judicial review under state constitutions when regulating federal elections, though federal courts retain a supervisory duty to ensure state courts do not transgress the ordinary bounds of judicial review and arrogate legislative power.

\section{Detailed Case Facts}
Following the 2020 decennial census, North Carolina's population grew by nearly one million residents, entitling the State to an additional seat in the U.S.\ House of Representatives (\href{https://www.law.cornell.edu/supremecourt/text/21-1271}{Cornell LII}; \href{https://www.courtlistener.com/opinion/10049660/moore-v-harper/}{CourtListener}). In November 2021, the Republican-controlled North Carolina General Assembly enacted three new redistricting maps along party lines: one for the State Senate, one for the State House, and one for North Carolina's federal congressional districts (N.C.\ Gen.\ Stat.\ Ann.\ \S\S~120-1, 120-2, 163-201).

Shortly thereafter, several groups of plaintiffs---including individual voters such as Rebecca Harper, the North Carolina League of Conservation Voters, and Common Cause---filed consolidated lawsuits in North Carolina state court against state legislative leaders (including Speaker of the North Carolina House of Representatives Timothy K.\ Moore) and state election officials (\href{https://www.law.cornell.edu/supremecourt/text/21-1271}{Cornell LII}). Plaintiffs alleged that the 2021 maps were extreme, intentional partisan gerrymanders violating multiple provisions of the North Carolina Constitution, including the Free Elections Clause (``[a]ll elections shall be free,'' Art.\ I, \S~10), the Equal Protection Clause, the Free Speech Clause, and the Free Assembly Clause.

In defense of the congressional map, the legislative defendants argued both (1) that partisan gerrymandering claims were nonjusticiable political questions under the North Carolina Constitution, and (2) under the ``independent state legislature theory'' (``ISLT''), that the Federal Elections Clause---which provides that ``[t]he Times, Places and Manner of holding Elections for Senators and Representatives, shall be prescribed in each State by the Legislature thereof,'' U.S.\ Const.\ art.\ I, \S~4, cl.\ 1---vests exclusive, unchecked authority in state legislatures to draw congressional maps free from state constitutional limits and state judicial review.

\section{Procedural History}
A three-judge panel of the Wake County Superior Court held a trial and found as a factual matter that the 2021 congressional map was ``a partisan outlier intentionally and carefully designed to maximize Republican advantage,'' yet entered judgment dismissing plaintiffs' claims on January 11, 2022, holding that partisan gerrymandering presented a nonjusticiable political question under state law (\href{https://www.law.cornell.edu/supremecourt/text/21-1271}{Cornell LII}).

In February 2022, the North Carolina Supreme Court (then holding a 4--3 Democratic majority) reversed in \textit{Harper v. Hall}, 380 N.C.\ 317, 868 S.E.2d 499 (2022) (\textit{Harper I}). Distinguishing \textit{Rucho v. Common Cause}, 588 U.S.\ 684 (2019), as a limit solely on federal courts, the court held that partisan gerrymandering claims were justiciable under the North Carolina Constitution, struck down the 2021 maps, rejected defendants' Elections Clause defense, enjoined use of the 2021 plans, and remanded for remedial map-drawing. When the General Assembly enacted a remedial congressional map on February 17, 2022 (which included a statutory ``trigger'' provision stating that the 2021 map would again become effective if the U.S.\ Supreme Court reversed \textit{Harper I}), the trial court rejected the remedial map and imposed an interim map drafted by Special Masters for the 2022 elections. The U.S.\ Supreme Court denied an emergency stay in March 2022 (over dissents by Justices Alito, Thomas, and Gorsuch) but granted certiorari on June 30, 2022 (\href{https://www.supremecourt.gov/opinions/22pdf/21-1271_3f14.pdf}{SCOTUS Slip Opinion}).

Oral argument took place on December 7, 2022. Nine days later, in \textit{Harper v. Hall}, 383 N.C.\ 89, 881 S.E.2d 156 (2022) (\textit{Harper II}), the outgoing North Carolina Supreme Court affirmed the trial court's rejection of the legislature's remedial congressional map. However, following the November 2022 elections, Republicans gained a 5--2 majority on the North Carolina Supreme Court. On January 20, 2023, legislative defendants petitioned for rehearing of \textit{Harper II}, asking the court to withdraw \textit{Harper II} and overrule \textit{Harper I} while expressly conceding that overruling \textit{Harper I} would not vacate its earlier judgment enjoining the 2021 plans. After granting rehearing, the North Carolina Supreme Court issued \textit{Harper v. Hall}, 886 S.E.2d 393 (N.C.\ 2023) (\textit{Harper III}) on April 28, 2023, overruling \textit{Harper I}, withdrawing \textit{Harper II}, holding partisan gerrymandering nonjusticiable under state law, and dismissing plaintiffs' claims with prejudice without reinstating the 2021 maps (\href{https://harvardlawreview.org/print/vol-137/moore-v-harper/}{Harvard Law Review}). Following two rounds of supplemental briefing on mootness and jurisdiction, the U.S.\ Supreme Court decided \textit{Moore v. Harper} on June 27, 2023.

\section{Judicial Votes}
\textbf{6--3 Decision Affirming the Judgment of the North Carolina Supreme Court in \textit{Harper I}} (\href{https://www.law.cornell.edu/supremecourt/text/21-1271}{Cornell LII}):
\begin{itemize}
\item \textbf{Majority (6):} Chief Justice John G.\ Roberts, Jr.\ delivered the opinion of the Court, joined in full by Justices Sonia Sotomayor, Elena Kagan, Brett M.\ Kavanaugh, Amy Coney Barrett, and Ketanji Brown Jackson.
\item \textbf{Concurrence (1):} Justice Brett M.\ Kavanaugh filed a concurring opinion while joining the Court's opinion in full.
\item \textbf{Dissent (3):} Justice Clarence Thomas filed a dissenting opinion, joined in full (Parts I, II, and III) by Justice Neil M.\ Gorsuch, and joined as to Part I (mootness and Article III jurisdiction) only by Justice Samuel A.\ Alito, Jr.
\end{itemize}

\section{Holding}
\textbf{Judgment:} The judgment of the North Carolina Supreme Court in \textit{Harper I} (380 N.C.\ 317, 868 S.E.2d 499) was \textbf{affirmed} (\href{https://www.law.cornell.edu/supremecourt/text/21-1271}{Cornell LII}).

\textbf{Rule Announced:} (1) \textit{Jurisdiction:} The case was neither moot under Article III nor non-final under 28 U.S.C.\ \S~1257(a) (applying \textit{Cox Broadcasting Corp. v. Cohn}, 420 U.S.\ 469 (1975)) because \textit{Harper III} overruled \textit{Harper I}'s reasoning without vacating its final judgment enjoining the 2021 maps, and state law provided that those maps would take effect again if the U.S.\ Supreme Court reversed. (2) \textit{Merits:} The Federal Elections Clause (U.S.\ Const.\ art.\ I, \S~4, cl.\ 1) does not vest exclusive, independent authority in state legislatures free from state constitutional constraints or the ordinary exercise of state judicial review. (3) \textit{Federal Backstop:} State courts nevertheless ``do not have free rein''; federal courts must ensure that state courts ``do not transgress the ordinary bounds of judicial review such that they arrogate to themselves the power vested in state legislatures to regulate federal elections.''

\textbf{Holding vs.\ Dicta:} The Court's core holdings are (1) that Article III and \S~1257(a) jurisdiction remained intact, (2) that the Elections Clause rejects the maximalist ISLT and subjects state legislatures to ordinary state constitutional constraints and state judicial review (with no procedure-versus-substance carve-out), and (3) that federal courts retain an outer-bounds check against state courts abrogating legislative authority. Because the Court explicitly declined to adopt a specific legal test for measuring when a state court transgresses ``ordinary bounds'' and declined to apply that principle to \textit{Harper I} (finding the issue forfeited by petitioners), its discussion of analogies to the Takings Clause, Contracts Clause, adequate-and-independent-state-grounds doctrine, and the separate writings in \textit{Bush v. Gore}, 531 U.S.\ 98 (2000), operates largely as guiding dicta---what commentators have dubbed ``ISLT-lite'' (\href{https://journals.law.harvard.edu/jol/wp-content/uploads/sites/86/2025/10/Morgan_JOL_Post-Galleys-2.pdf}{Harvard J.\ on Legis.}).

\section{Analysis of Opinions}
\subsection{Majority Opinion (Chief Justice Roberts)}
Chief Justice Roberts divided the Court's opinion into four main parts (\href{https://www.law.cornell.edu/supremecourt/text/21-1271}{Cornell LII}).

\textbf{1. Jurisdiction and Mootness (Part II).} Roberts held that the dispute remained live under Article III because \textit{Harper I}'s mandate had already issued, legislative defendants conceded they remained bound by its injunction against the 2021 maps, and North Carolina's 2022 statutory trigger clause (like the statute in \textit{Hunt v. Cromartie}, 526 U.S.\ 541 (1999)) would automatically reinstate the 2021 maps upon reversal. Under 28 U.S.C.\ \S~1257(a), \textit{Harper I} satisfied the second category of \textit{Cox Broadcasting} because the Elections Clause issue was finally decided by the state's highest court and survived subsequent state proceedings.

\textbf{2. Historical Roots of Judicial Review (Part III).} Tracing judicial review prior to \textit{Marbury v. Madison}, 5 U.S.\ (1 Cranch) 137 (1803), Roberts cited Founding-era state decisions including \textit{Trevett v. Weeden} (R.I.\ 1786), North Carolina's own \textit{Bayard v. Singleton}, 1 Mort.\ 48 (1787), and \textit{Holmes v. Walton} (N.J.\ 1780), alongside remarks by James Madison, Elbridge Gerry, and Alexander Hamilton (\textit{The Federalist Nos.\ 78 and 81}). Because state judicial review was already a ``fundamental principle'' at the Founding, the Elections Clause did not silently carve out an exception to it.

\textbf{3. Precedent, Constitutional Structure, and Founding Practice (Part IV).} Roberts relied on a century of precedent---\textit{Ohio ex rel.\ Davis v. Hildebrant}, 241 U.S.\ 565 (1916) (popular referendum), \textit{Smiley v. Holm}, 285 U.S.\ 355 (1932) (gubernatorial veto), and \textit{Arizona State Legislature v. Arizona Independent Redistricting Commission}, 576 U.S.\ 787 (2015) (independent commission)---to establish that prescribing election rules is a lawmaking function subject to state constitutional constraints. Legislatures are ``creatures of the State Constitutions'' operating under a two-fold commission: both state and federal constitutions restrain them. Roberts distinguished \textit{McPherson v. Blacker}, 146 U.S.\ 1 (1892), as involving no state-constitutional conflict, and \textit{Leser v. Garnett}, 258 U.S.\ 130 (1922), and \textit{Hawke v. Smith}, 253 U.S.\ 221 (1920), as involving Article V ratification rather than lawmaking. He rejected petitioners' attempt to limit \textit{Smiley} to ``procedural'' rather than ``substantive'' state constitutional limits, noting that precedent drew no such line and that the distinction between procedure and substance is ``hazy.'' Finally, Founding-era state constitutions (including Delaware's 1792 Constitution, Maryland's 1810 amendment, numerous ``by ballot'' clauses, and state regulations of delegates under Article V of the Articles of Confederation) confirmed that state constitutions routinely regulated federal elections.

\textbf{4. Outer Federal Limits on State Courts (Part V).} Even while rejecting maximalist ISLT, Roberts emphasized that ``state courts do not have free rein.'' Drawing analogies to the Takings Clause (\textit{Tyler v. Hennepin County}, 598 U.S.\ 631 (2023); \textit{Phillips v. Washington Legal Foundation}, 524 U.S.\ 156 (1998)), the Contracts Clause (\textit{Indiana ex rel.\ Anderson v. Brand}, 303 U.S.\ 95 (1938)), adequate-and-independent-state-grounds cases (\textit{NAACP v. Alabama ex rel.\ Patterson}, 357 U.S.\ 449 (1958)), and the separate opinions of Chief Justice Rehnquist and Justice Souter in \textit{Bush v. Gore}, 531 U.S.\ 98 (2000), Roberts held that federal courts must ensure state courts do not ``transgress the ordinary bounds of judicial review'' and arrogate legislative authority. Because petitioners expressly disclaimed any argument that the North Carolina Supreme Court misread state law, the Court adopted no specific test and affirmed \textit{Harper I}.

\subsection{Concurring Opinion (Justice Kavanaugh)}
Joining the majority in full, Justice Kavanaugh wrote separately to propose a concrete standard for federal review of state-court election rulings (\href{https://www.law.cornell.edu/supremecourt/text/21-1271}{Cornell LII}). Comparing Chief Justice Rehnquist's concurrence in \textit{Bush v. Gore} (whether the state court ``impermissibly distorted'' state law ``beyond what a fair reading required''), Justice Souter's dissent in \textit{Bush v. Gore} (whether it exceeded ``the limits of reasonable'' interpretation), and the Solicitor General's amicus proposal (``truly aberrant''), Kavanaugh argued that all three capture the same principle: federal review ``should be deferential, but deference is not abdication.'' He endorsed Rehnquist's formulation, argued it should apply to both state statutory and state constitutional interpretations, and stressed that federal courts must measure a state court's ruling against the state law as it existed prior to that court's decision.

\subsection{Dissenting Opinion (Justice Thomas, joined by Justice Gorsuch, and by Justice Alito as to Part I)}
Justice Thomas dissented on three grounds (\href{https://www.law.cornell.edu/supremecourt/text/21-1271}{Cornell LII}).

\textbf{1. Mootness and Advisory Opinions (Part I, joined by Gorsuch and Alito).} Thomas argued the case was ``indisputably moot'' and the majority's opinion ``plainly advisory.'' Once \textit{Harper III} overruled \textit{Harper I}, affirmed the original January 2022 trial-court judgment, and dismissed plaintiffs' claims with prejudice on state-law nonjusticiability grounds, petitioners had won complete relief; the interlocutory injunction in \textit{Harper I} merged into the final judgment. Courts adjudicate parties' rights, not the abstract validity of statutes, and a state legislature's statutory trigger clause cannot manufacture an Article III controversy.

\textbf{2. Substantive vs.\ Procedural ISLT (Part II, joined by Gorsuch only).} On the merits, Thomas defended a narrower version of ISLT rooted in three premises: (1) a single State's people cannot limit powers delegated by the Federal Constitution (\textit{McCulloch v. Maryland}; \textit{Leser v. Garnett}); (2) regulating federal elections is a federally delegated power, not a reserved state power (\textit{Cook v. Gralike}; \textit{U.S.\ Term Limits, Inc. v. Thornton}); and (3) ``the Legislature'' means the state's lawmaking power as constituted by its state constitution. Thus, while state constitutions may define \textit{who} makes laws and \textit{how} (procedural rules like gubernatorial vetoes in \textit{Smiley}), they cannot impose \textit{substantive} limits on \textit{what} federal election rules the legislature may enact.

\textbf{3. Unworkability of the Majority's Federal-Review Standard (Part III, joined by Gorsuch only).} Thomas warned that extending \textit{Bush v. Gore}-style review from state statutes to broad, open-textured state constitutions would federalize state constitutional law without intelligible standards, swell federal dockets, and risk having federal courts decide election winners through expedited, haphazard judgments about whether a state court exceeded ``the bounds of ordinary judicial review.''

\section{Examples: Future Applications}
\begin{enumerate}
\item \textbf{Same-Side Application (State Court Enforcing Express State Constitutional Redistricting Rules):} A state constitution expressly mandates that congressional districts be contiguous, compact, and preserve county boundaries where practicable, and a state supreme court invalidates a legislature's map for violating those textual criteria. Under \textit{Moore}, the state court is squarely exercising ordinary judicial review under substantive state constitutional limits, and the Elections Clause provides no basis for federal reversal.

\item \textbf{Same-Side Application (State Court Enforcing Absentee-Voting Protections Under Settled State Precedent):} Relying on decades of state supreme court precedent interpreting the state constitution's ``free and equal elections'' clause, a state court strikes down a newly enacted statute restricting mail-in voting drop boxes for both state and federal elections. Because the court applies established state interpretive methods and precedent, it operates within the ordinary bounds of judicial review under \textit{Moore}.

\item \textbf{Opposite-Side Application (State Court Rewriting Unambiguous Statutory Deadlines Without Legal Basis):} Days before a federal election, a state trial court extends a clear, unambiguous statutory ballot-receipt deadline by two weeks based solely on generalized ``policy fairness,'' without finding any constitutional violation or grounding the remedy in any state statute or equitable precedent. Under \textit{Moore}'s outer-bounds limit (and Kavanaugh's concurrence), a federal court could intervene because the state court transgressed ordinary judicial review and arrogated legislative power under the Elections Clause.

\item \textbf{Opposite-Side Application (State Court Imposing Its Own Congressional Map Without Finding the Enacted Map Unlawful):} A state court enjoins a legislature's duly enacted congressional map and imposes a map drawn by a plaintiff's expert without ever holding that the legislative map violated any state or federal constitutional provision. Replacing legislative districts without a predicate finding of illegality exceeds ordinary judicial review and intrudes on the authority reserved to ``the Legislature'' under Article I, \S~4.

\item \textbf{Fence-Sitter (Novel State Constitutional Interpretation During an Election Emergency):} During a severe hurricane or pandemic, a state supreme court invokes a broadly worded state constitutional ``free elections'' guarantee for the first time in state history to extend a statutory mail-in ballot deadline by three days solely for that election. This is a classic fence-sitter: defenders argue that giving meaning to broad state constitutional text in unforeseen crises is quintessential common-law judging (\href{https://electionlawblog.org/2023/carolyn-shapiro-moore-v-harper-and-state-courts/}{Shapiro, Election Law Blog}), whereas challengers invoking Kavanaugh's pre-existing-law baseline argue that a novel judicial alteration of unambiguous legislative deadlines crosses the line into lawmaking.
\end{enumerate}

\section{Critique}
\textit{Moore v. Harper} drew widespread attention as a landmark separation-of-powers and federalism ruling, praised for averting a destabilizing maximalist ISLT yet criticized from both right and left for leaving an undefined federal judicial backstop.

\textbf{Originalist and Textualist Perspectives:} Originalist scholars were sharply divided on the underlying Elections Clause question, though a broad coalition of prominent originalists---including Akhil Reed Amar, Vikram David Amar, and Steven Calabresi---celebrated the rejection of maximalist ISLT. As Vikram Amar argued, ``proper originalism'' looking to public meaning, the Articles of Confederation's parallel delegate-selection clause, and uniform early post-ratification state constitutions (all of which regulated federal elections without objection from George Washington or the First Congress) decisively refuted ISLT, whereas petitioners' reliance on the discredited ``Pinckney Plan'' failed historical scrutiny (\href{https://verdict.justia.com/2022/11/09/what-the-isl-moore-v-harper-case-can-tell-us-about-principled-originalism}{Amar, Justia Verdict}; \href{https://www.brennancenter.org/our-work/court-cases/annotated-guide-amicus-briefs-moore-v-harper}{Brennan Center Amicus Guide}). Other conservative and originalist scholars, such as William Baude and Michael McConnell, supported a narrower middle ground---affirming state judicial review while questioning whether state courts possess independent authority to draw replacement congressional maps rather than merely invalidating unlawful ones (\href{https://redistricting.lls.edu/wp-content/uploads/UT-gardner-20260214-mtd-opp.pdf}{Baude \& McConnell cited in Utah Redistricting Litigation}). Meanwhile, formalist critics echoed Justice Thomas's dissent on Article III, arguing that the majority stretched \textit{Cox Broadcasting} and mootness doctrine to issue an advisory opinion after \textit{Harper III} had already ended the North Carolina litigation (\href{https://harvardlawreview.org/print/vol-137/moore-v-harper/}{Harvard Law Review}).

\textbf{Progressive and Doctrinal Perspectives:} While voting-rights advocates welcomed the defeat of maximalist ISLT, progressive and election-law scholars---including Richard Hasen, Richard Pildes, Carolyn Shapiro, Leah Litman, and Katherine Shaw---warned that Part V planted a doctrinal ``time bomb'' (\href{https://harvardlawreview.org/print/vol-137/moore-v-harper/}{Harvard Law Review}; \href{https://nyulawreview.org/wp-content/uploads/2024/12/99-NYU-L-Rev-2049.pdf}{NYU Law Review}). Litman and Shaw cautioned that the undefined ``ordinary bounds of judicial review'' standard (and especially Kavanaugh's ``anti-novelty'' gloss) invites federal judges to impose their own preferred textualist methodology on state supreme courts, chilling state-level development of voting rights (\href{https://repository.law.umich.edu/cgi/viewcontent.cgi?article=3950&context=facarticles}{Litman \& Shaw, Mich.\ Law}). Similarly, the \href{https://harvardlawreview.org/print/vol-137/moore-v-harper/}{\textit{Harvard Law Review}} noted the irony that \textit{Moore} curbed potential legislative mischief only by elevating an unelected federal judiciary as the ultimate supervisor of fifty states' election jurisprudence, reviving the controversial \textit{Bush v. Gore} framework in a Supreme Court majority opinion for the first time. Shapiro argued that federal oversight, if any, should rest on narrow due-process reliance principles rather than an indeterminate Elections Clause baseline (\href{https://electionlawblog.org/2023/carolyn-shapiro-moore-v-harper-and-state-courts/}{Shapiro, Election Law Blog}).

\textbf{Underlying Competing Values:} \textit{Moore v. Harper} reflects a multi-cornered institutional dilemma---a ``new countermajoritarian difficulty'' (\href{https://harvardlawreview.org/print/vol-137/moore-v-harper/}{Harvard Law Review})---pitting \textit{state legislative accountability and federal textual fidelity} (respecting Article I's assignment of election rules to ``the Legislature'') against \textit{intragovernmental checks and balances} (preventing gerrymandered legislatures from rigging rules unchecked by state constitutions) and \textit{judicial federalism} (preserving state supreme courts as final arbiters of state law without inviting federalized election litigation on the emergency docket).

\section{Key Quotations}
\begin{quote}
``The Elections Clause does not insulate state legislatures from the ordinary exercise of state judicial review.''
\end{quote}
\noindent --- Chief Justice John G.\ Roberts, Jr., Majority Opinion, \textit{Moore v. Harper}, 600 U.S.\ 1 (2023) (\href{https://www.law.cornell.edu/supremecourt/text/21-1271}{Cornell LII}; \href{https://www.courtlistener.com/opinion/10049660/moore-v-harper/}{CourtListener}).

\begin{quote}
``The legislature acts both as a lawmaking body created and bound by its state constitution, and as the entity assigned particular authority by the Federal Constitution. Both constitutions restrain the legislature's exercise of power.''
\end{quote}
\noindent --- Chief Justice John G.\ Roberts, Jr., Majority Opinion (\href{https://www.law.cornell.edu/supremecourt/text/21-1271}{Cornell LII}).

\begin{quote}
``We do not adopt these or any other test by which we can measure state court interpretations of state law in cases implicating the Elections Clause. The questions presented in this area are complex and context specific. We hold only that state courts may not transgress the ordinary bounds of judicial review such that they arrogate to themselves the power vested in state legislatures to regulate federal elections.''
\end{quote}
\noindent --- Chief Justice John G.\ Roberts, Jr., Majority Opinion (\href{https://www.law.cornell.edu/supremecourt/text/21-1271}{Cornell LII}).

\begin{quote}
``Federal court review of a state court's interpretation of state law in a federal election case should be deferential, but deference is not abdication.''
\end{quote}
\noindent --- Justice Brett M.\ Kavanaugh, Concurring Opinion (\href{https://www.law.cornell.edu/supremecourt/text/21-1271}{Cornell LII}).

\begin{quote}
``In the end, I fear that this framework will have the effect of investing potentially large swaths of state constitutional law with the character of a federal question not amenable to meaningful or principled adjudication by federal courts.''
\end{quote}
\noindent --- Justice Clarence Thomas, Dissenting Opinion (\href{https://www.law.cornell.edu/supremecourt/text/21-1271}{Cornell LII}).

\end{document}
